%================================================================
\section{Foil 2: the estimand and the calibration check}\label{app:foil2}

\paragraph{The H\'ajek target and the kernel truth} The H\'ajek mean in
Equation~\eqref{eq:hajek} targets the inverse-probability exposure contrast. The
kernel truth $\dspill$ in Equation~\eqref{eq:spill} averages the unweighted
realized contrasts. The exposure probabilities vary widely across people in the
two unseeded strata, so the two functionals weight people differently. On this run
the two are close enough that no containment flag depends on the difference: in the
panels we examined, outcomes hardly track the inverse-probability weights. That
comparison was made offline on paired panels and is not a committed result. It is
a fact about this kernel, graph and parameter point; we did not run another kernel.
The exposure map uses neighbors' seeding, whereas the kernel's peer term uses
neighbors' outcomes and the effect travels more than one step, so the map is not
exact for the kernel. Foil~2 cannot target $\theta_2$ directly: the policy arm has
no lottery, so given the graph each person's assignment probability is $0$ or $1$
and inverse-probability weights are not defined.

\paragraph{The pre-registered calibration check} The script
\texttt{scripts/check\_foil2\_coverage.py} is a required check of Foil~2's
interval. Its gate uses outcomes that obey the exposure map exactly, with a true
spillover of $\GateTruth$, and it draws $\GateAssignments$ pilot assignments. The
gate rule was written before the run: the gate fails if the ratio of the mean
estimated variance to the empirical variance, plus $1.645$ bootstrap standard
errors, is below one. The gate passes, with a mean estimate of $\GateMeanEstimate$
and a ratio of $\GateRatio$ (bootstrap SE $\GateRatioBootSE$). The result is:
\begin{quote}
In a fixed-world repeated-assignment check whose outcomes obey the exposure map
exactly, the nominal $\ConfidencePct\%$ Foil~2 Wald interval covered
$\GateCoveragePct\%$ of assignments (exact $95\%$ binomial interval
$\GateCoverageLoPct\%$--$\GateCoverageHiPct\%$). The average estimated variance
was $\GateRatio$ times the empirical design variance, so the shortfall is not
explained by average variance underestimation. It reflects the finite-sample
behavior of this studentized normal interval, with the implemented H\'ajek
linearization and simulated joint exposure probabilities.
\end{quote}
The check is specific to this implementation and is not a test of the theorem of
\citet{aronow2017}, whose exact result is for Horvitz--Thompson totals with exact
joint probabilities. The code uses H\'ajek means, a first-order linearization and
one simulated table of joint probabilities.

The script also reports the kernel itself, where the exposure map is not exact.
Over $\CheckPanels$ kernel panels, the nominal $\ConfidencePct\%$ interval
contains the truth at a rate of $\KernelCovStar$ at $\cstar$, $\KernelCovMinI$ at
the min-I point and $\KernelCovMedI$ at the median-I point. Redrawing only the
assignment while holding the noise fixed gives $\DesignCovStar$, $\DesignCovMinI$
and $\DesignCovMedI$. The code reports these numbers and does not gate on them.

The spillover contrast is the case in which the Young bound matters least. In the
independent review of amendment A2, the Young share of the spillover contrast's
variance on the $\cstar$ panel was small, \MISSING{YoungShareSpillStar}, whereas
the A1 and A2 contrasts include the $\Sb$ stratum and carry a much larger Young
share (Section~\ref{sec:a2results}). This is consistent with Foil~2's interval
covering a little under nominal in its check while the A1 and A2 intervals are
conservative.

